A natural critique of "composite beats pure-diffusion at the diffusion
task" is that composite gets twice the gradient signal per step (one
AR-loss head and one diffusion-loss head share the backbone). If
pure-diffusion were trained for $2\times$ the steps, would the gap
disappear? This section shows it does not.

\subsection{Setup}

We compare three conditions at $200$M parameters:

\begin{itemize}
  \item Composite, $3{,}000$ steps, diffusion-axis NLL on held-out chunk
  \item Pure-diffusion, $3{,}000$ steps, same chunk
  \item \textbf{Pure-diffusion, $6{,}000$ steps} (the compute-matched
        control), same chunk
\end{itemize}

\diffNLL{} measurement: mask $50\%$ of answer-region tokens uniformly
at random and compute cross-entropy on masked positions under
bidirectional forward.

\subsection{Results}

\begin{table}[h]
\centering
\begin{tabular}{lrrr}
\toprule
Model & Training steps & Diff loss share & \diffNLL{} \\
\midrule
Composite ($200$M) & $3{,}000$ & $\sim 25\%$ & \textbf{5.42} \\
Pure-diffusion ($200$M) & $3{,}000$ & $100\%$ & 6.13 \\
\textbf{Pure-diffusion ($200$M)} & $\mathbf{6{,}000}$ & $\mathbf{100\%}$ & \textbf{6.11} \\
\bottomrule
\end{tabular}
\caption{Diffusion-axis NLL on held-out GSM8K answer tokens, $50\%$
mask ratio. Composite wins by $-0.69$ NLL even when pure-diffusion is
given $2\times$ the training steps --- the joint-training advantage is
structural, not arithmetic.}
\label{tab:compute-matched}
\end{table}

Pure-diffusion at $6{,}000$ steps barely improves over $3{,}000$
steps ($6.11$ vs $6.13$). Composite at $3{,}000$ steps still beats
both ($5.42$). \textbf{The composite advantage cannot be explained as
"composite saw more gradient signal."}

\subsection{Multi-Scale Replication}

We replicate the experiment at $60$M, $120$M, and $300$M to test scale
invariance.
\textsc{[TODO: fill from \texttt{e5/results/e3b\_multiscale\_d2/summary.json}]}.

\subsection{Why this matters}

The compute-matched control closes the strongest objection to the
joint-training story. The remaining alternative explanations are:

\begin{enumerate}
  \item \emph{Bidirectional attention itself helps mask-fill, and
        the AR loss component does not actively hurt.} This is
        consistent with our data and is essentially the
        joint-training-as-architecture-prior story.
  \item \emph{The AR loss term constrains the backbone in a way that
        happens to also be helpful for diffusion at this scale.}
        Plausible, scale-dependent, and would predict the advantage
        shrinks at much larger scales (we cannot test).
\end{enumerate}

Both readings are favourable to composite design. The reading we
\textbf{rule out} is "composite wins because it gets more total
training compute" --- the compute-matched control falsifies that.
